\documentclass[11pt]{article}
\usepackage[margin=23mm]{geometry}
\usepackage{amsmath,amssymb,amsthm}
\usepackage[T1]{fontenc}
\usepackage{lmodern}
\usepackage{tikz-cd}
\usepackage[colorlinks=true,urlcolor=blue,linkcolor=blue,hypertexnames=false]{hyperref}
\newcommand{\I}{\mathbb I}
\newcommand{\C}{\mathbf C}
\newcommand{\Q}{\mathbb Q}
\newcommand{\Z}{\mathbb Z}
\newcommand{\Nm}{\operatorname{Nm}}
\newcommand{\Gal}{\operatorname{Gal}}
\newcommand{\Hom}{\operatorname{Hom}}
\newcommand{\Br}{\operatorname{Br}}
\newcommand{\inv}{\operatorname{inv}}
\newtheorem*{target}{Theorem VII.8.1}
\newenvironment{authorresult}[2]{\par\medskip\noindent\textbf{#1 VII.#2.}\ \itshape}{\par\medskip}
\title{Global reciprocity on principal id\`eles\\and the sum of local invariants}
\author{J. S. Milne\\\small Source: \emph{Class Field Theory}, v4.03, 6 August 2020}
\date{}
\begin{document}
\maketitle
\noindent\small
Private reading transcription of the author's Theorem VII.8.1, its supplied
cyclotomic proof, and the necessary Brauer-group support in VII\S7.
The selected core occupies physical PDF pages 228--232 (printed pages
219--223). The local cyclotomic formulas used in Example VII.8.2 are retained
from I, Example 3.13, physical page 53 (printed page 44), with the separate
published correction identified below. The complete fixed original is retained
at \url{https://www.jmilne.org/math/CourseNotes/CFT.pdf}.
Mathematical expressions and the author's argument are retained, with
connective prose and layout normalized. This is published-PDF transcription,
not author-native TeX or a new proof.
\normalsize

% BEGIN CONTEXT VII.8.1 CONVENTIONS
\section*{Notation and preceding source results}
For a number field $K$, $\I_K$ is its id\`ele group and
$\C_K=\I_K/K^\times$. Write
\[
 H^2(L/K)=H^2(\Gal(L/K),L^\times),\qquad
 H^2(/K)=H^2(\Gal(K^{\mathrm{al}}/K),(K^{\mathrm{al}})^\times).
\]
As explained in Chapter IV, these identify with $\Br(L/K)$ and $\Br(K)$.
For a place $v$ of $K$, $L^v$ denotes the completion at a place above $v$;
$G^v$ is the decomposition group. The local image of a cohomology class
$\alpha$ is denoted $\alpha_v$. The source abbreviates
$\inv_v(\alpha_v)$ to $\inv_v(\alpha)$.
The final sentence of VII.8.6 writes $H^2(K)$ after earlier using
$H^2(/K)$ for its ambient absolute class; that printed notation is retained.

For a finite abelian extension $L/K$ with Galois group $G$, the author has
defined the homomorphism
\[
 \phi_{L/K}:\I_K\longrightarrow G,\qquad
 \phi_{L/K}(\mathbf a)=\prod_v\phi_v(a_v)
 \quad\bigl(\mathbf a=(a_v)\bigr),
\]
from the local Artin maps. The background used here includes local class
field theory and its invariant map, the decomposition of id\`ele cohomology
(VII.2.5), cup products, Hilbert 90, and inflation--restriction.
The preceding second inequality (VII.5.1) gives $H^1(G,\C_L)=0$ and
\[
 (\I_K:K^\times\Nm_{L/K}\I_L)\mid[L:K].
\]
Its complete algebraic proof is a separate, linked reader in this same
source work; that core is not counted again here.
In the proof of VII.8.5 the author writes the abelian target group
additively; there $\phi_{L/K}(a)=0$ denotes its identity element.
% END CONTEXT VII.8.1 CONVENTIONS

% BEGIN PRIMARY STATEMENT VII.8.1
\section*{The target statement}
\begin{target}
\begin{enumerate}
\item[(a)] Let $L/K$ be a finite abelian extension of number fields. Then
$\phi_{L/K}$ takes the value $1$ on the principal id\`eles
$K^\times\subset\I_K$.
\item[(b)] Let $L/K$ be a finite Galois extension of number fields. Then
\[
 \sum_v\inv_v(\alpha)=0\qquad\text{for every }\alpha\in H^2(L/K).
\]
\end{enumerate}
\end{target}
% END PRIMARY STATEMENT VII.8.1

% BEGIN PRIMARY PROOF VII.8.1 BRAUER INJECTIVITY
\section*{The author's Brauer-group support}
\begin{authorresult}{Theorem}{7.1}
For every Galois extension $L/K$ of number fields, possibly infinite, the
canonical map
\[
 H^2(L/K)\longrightarrow\bigoplus_v H^2(L^v/K_v)
\]
is injective.
\end{authorresult}
\begin{proof}
First suppose that $L/K$ is a finite Galois extension with Galois group $G$.
Because $H^1(G,\C_L)=0$, the cohomology sequence of
\[
 0\longrightarrow L^\times\longrightarrow\I_L\longrightarrow\C_L
 \longrightarrow0
\]
gives
\[
 0\longrightarrow H^2(G,L^\times)\longrightarrow H^2(G,\I_L)
 \longrightarrow\cdots.
\]
But $H^2(G,L^\times)=H^2(L/K)$ and, by VII.2.5,
\[
 H^2(G,\I_L)
 =\bigoplus_v H^2(G^v,L^{v\times})
 =\bigoplus_v H^2(L^v/K_v),
\]
which proves the finite case. For an infinite extension, pass to the limit
over its finite Galois subextensions.
\end{proof}
% END PRIMARY PROOF VII.8.1 BRAUER INJECTIVITY

% BEGIN PRIMARY PROOF VII.8.1 CYCLOTOMIC SPLITTING
An extension $L/K$ is called \emph{cyclotomic} if $L\subset K[\zeta]$ for
some root of $1$, $\zeta$.
\begin{authorresult}{Proposition}{7.2}
For every $\beta\in H^2(/K)$, there is a cyclic cyclotomic extension $L$ of
$K$ such that $\beta$ maps to zero in $H^2(/L)$.
\end{authorresult}
\begin{proof}
Theorem VII.7.1 shows that $\beta$ is determined by its images in $H^2(K_v)$,
and hence by $\inv_v(\beta_v)\in\Q/\Z$ (III, Theorem 2.1).
For every finite extension $L/K$ and prime $w$ of $L$ above $v$,
\[
 \inv_w(\beta|_L)=[L_w:K_v]\inv_v(\beta).
\]
Thus we must find a cyclic cyclotomic $L/K$ such that
\[
 [L^v:K_v]\inv_v(\beta_v)=0\pmod{\Z}\quad\text{for all }v.
\]
Because the map from $H^2(L/K)$ into $H^2(/K)$ maps into the direct sum of
the local groups, $\inv_v(\beta_v)=0$ for almost all $v$.
Consequently some integer $m>0$ satisfies
$m\inv_v(\beta_v)=0$ for every $v$. The next lemma supplies $L$.
\end{proof}

\begin{authorresult}{Lemma}{7.3}
Given a number field $K$, a finite set $S$ of finite primes of $K$, and an
integer $m>0$, there is a totally complex cyclic cyclotomic extension $L$
of $K$ such that $m\mid[L^v:K_v]$ for every $v\in S$.
\end{authorresult}
\begin{proof}
It suffices to prove the assertion for $\Q$ and $m[K:\Q]$, so assume $K=\Q$.
Let $\ell$ be prime and let $\zeta$ be a primitive $\ell^r$th root of $1$
with $r>2$. Then
\[
 \Gal(\Q[\zeta]/\Q)\simeq(\Z/\ell^r\Z)^\times,
\]
and
\[
 (\Z/\ell^r\Z)^\times\simeq
 \begin{cases}
 \Delta\times C(\ell^{r-1})&\ell\text{ odd},\\
 \Delta\times C(2^{r-2})&\ell=2.
 \end{cases}
\]
Here $\Delta$ has order $\ell-1$ and $2$ in the two cases, and $C(t)$ is a
cyclic group of order $t$ (Serre 1970, II, 3.2).
Therefore $L(\ell^r)\overset{\mathrm{def}}=\Q[\zeta]^\Delta$ is a cyclic
cyclotomic extension of $\Q$ of degree $\ell^{r-2}$ or $\ell^{r-3}$.
\footnote{Source wording note: the degrees $\ell^{r-2}$ and $\ell^{r-3}$
are the exponents printed on physical page 229 (printed page 220).
This transcription retains that source sentence and supplies no replacement
degree formula or claimed author erratum.}

Consider $\Q_p[\zeta]$. If $p=\ell$, $\Q[\zeta]$ is totally ramified over
$p$, and
\[
 [\Q_p[\zeta]:\Q_p]=[\Q[\zeta]:\Q]=\varphi(\ell^r).
\]
If $p\ne\ell$, $\Q[\zeta]$ is totally unramified over $p$, and
$[\Q_p[\zeta]:\Q_p]=t$, where $t$ is the smallest integer such that
$\ell^r\mid p^t-1$. In either case, this local degree tends to infinity
as $r\to\infty$. Thus for every $p$, $[L(\ell^r)^p:\Q_p]$ is a power of
$\ell$ tending to infinity with $r$.

A product of cyclic groups of distinct prime-power orders is cyclic.
Therefore, for distinct primes $\ell_1,\ldots,\ell_s$,
\[
 L=L(\ell_1^{r_1})\cdots L(\ell_s^{r_s})
\]
is cyclic. By choosing $\ell_1^{r_1}\cdots\ell_s^{r_s}$ sufficiently large,
the local degrees $[L^p:\Q_p]$ are divisible by $m$ for all $p\in S$.
\end{proof}

The author expresses these results more concretely: a central simple
algebra over $K$ that splits over every $K_v$ splits over $K$; every central
simple algebra over $K$ splits over a cyclic cyclotomic extension of $K$.
% END PRIMARY PROOF VII.8.1 CYCLOTOMIC SPLITTING

% BEGIN CONTEXT VII.8.1 LOCAL CYCLOTOMIC FORMULAS
\section*{The local formulas used by the author}
The following are the portions of I, Example 3.13 used in VII, Example 8.2.
They describe the local Artin map $\phi_p:\Q_p^\times\to\Gal(L/\Q_p)$
when $L=\Q_p[\zeta]$ and $\zeta$ is a primitive $n$th root of $1$.

\noindent\textbf{I.3.13(a).}
If $n$ is prime to $p$, then $L/\Q_p$ is unramified, of degree equal to the
residue-field degree. The residue field is $\mathbb F_{p^f}$, where $p^f$
is the smallest power of $p$ with $n\mid p^f-1$.
The map sends $u p^t$ to the $t$th power of Frobenius, and its kernel is
$\Z_p^\times\langle p^f\rangle$.

\noindent\textbf{I.3.13(b), the used formula as printed.}
When $n=p^r$, write $a=u p^t$, and let $u_0\in\Z$ represent the class of
$u$ in $(\Z_p/p^r\Z_p)^\times$. The original page writes that
$\phi_p(a)$ sends $\zeta$ to $\zeta^{u^{-1}}$, and gives its kernel as
\[
 \{u p^m\mid u\equiv1\pmod{p^r},\ m\in\Z\}.
\]
Only the local-action formulas required for the selected proof are quoted
here; the complete example remains in the original fixed work.

\paragraph{Separate published correction.}
Under the CFT v4.03 section of the author's errata page, the typographical
note attributed to Raghuram Sundararajan says that the last line of
I.3.13(b) should use
\[
 \phi_p(a):\zeta\longmapsto\zeta^{u_0^{-1}}.
\]
The original and the correction are kept distinct.
The complete saved page is \url{https://www.jmilne.org/math/CourseNotes/errata.html}.
The correction to composite part (c) remains in that saved source; this
reader uses the prime-power route and does not transcribe (c).
% END CONTEXT VII.8.1 LOCAL CYCLOTOMIC FORMULAS

% BEGIN PRIMARY PROOF VII.8.1 RATIONAL CYCLOTOMIC CASE
\section*{The author's proof: the cyclotomic case}
\begin{authorresult}{Example}{8.2}
We verify part (a) for $\Q[\zeta_m]/\Q$, where $\zeta_m$ is a primitive
$m$th root of $1$.
\end{authorresult}
Identify $\Gal(\Q[\zeta_m]/\Q)$ with $(\Z/m\Z)^\times$.
Here $a\in\Q^\times$ is viewed as a principal id\`ele, as in part (a).
For $n$ relatively prime to $m$, $[n]$ denotes the automorphism sending
$\zeta_m$ to $\zeta_m^n$. It suffices to show
$\phi(a)|_{\Q[\zeta_{\ell^r}]}=1$ for all $\ell\mid m$.
We may assume $m=\ell^r$, $m\ne2$.

The homomorphism
\[
 \phi_\infty:\mathbb R^\times/\Nm(\mathbb C^\times)
 \longrightarrow\Gal(\Q[\zeta_m]/\Q)
\]
sends a negative real number to complex conjugation, so
$\phi_\infty(a)=[\operatorname{sign}(a)]$.
Write $a=u p^s\in\Q_p^\times$. If $p\ne\ell$, then $p$ is unramified in
$\Q[\zeta_m]$ and $\phi_p(a)$ acts as the $s$th power of Frobenius at $p$:
\[
 \phi_p(u p^s)=[p^s].
\]
The prime $\ell$ is totally ramified in $\Q[\zeta_m]$, and
\[
 \phi_\ell(a)=[u^{-1}]
\]
by I.3.13.
\footnote{The source records Catherine O'Neil's observation that this uses
I\S3 despite the earlier promise that I\S\S2--4 could be skipped.
The author mentions an elementary alternative and leaves its clarification
to a future version. The supplied local-Artin route is retained here.}

It suffices to show $\phi(a)=1$ in the three cases
$a=-1$, $a=\ell$, and $a=q$ for a prime $q\ne\ell$.
The author gives
\[
 \phi_p(-1)=
 \begin{cases}
 [-1]&p=\infty,\\
 [-1]&p=\ell,\\
 [1]&p\ne\ell,\infty,
 \end{cases}
 \qquad
 \phi_p(\ell)=
 \begin{cases}
 [1]&p=\ell,\\
 [1]&p\ne\ell,
 \end{cases}
\]
and
\[
 \phi_p(q)=
 \begin{cases}
 [q]&p=q,\\
 [q^{-1}]&p=\ell,\\
 [1]&p\ne\ell,q.
 \end{cases}
\]
In every case, $\prod_p\phi_p(a)=1$.
% END PRIMARY PROOF VII.8.1 RATIONAL CYCLOTOMIC CASE

% BEGIN PRIMARY PROOF VII.8.1 BASE CHANGE
\begin{authorresult}{Lemma}{8.4}
\begin{enumerate}
\item[(a)] If part (a) holds for $L/K$, it holds for every subextension.
\item[(b)] If part (a) holds for $L/K$, it holds for $LK'/K'$ for every
number field $K'\supset K$.
\end{enumerate}
\end{authorresult}
\begin{proof}
(a) Suppose $L\supset K'\supset K$. Then $\phi_{K'/K}$ is the composite of
$\phi_{L/K}$ and restriction $\Gal(L/K)\to\Gal(K'/K)$, because this is
true for the local Artin maps.

(b) Put $L'=LK'$. For each prime $w$ of $K'$, there is a commutative diagram
\[
\begin{tikzcd}[column sep=large,row sep=large]
K_w'{}^\times\arrow[r,"\phi_w"]\arrow[d,"\Nm"']&
\Gal(L'^w/K'_w)\arrow[d,hook]\\
K_v^\times\arrow[r,"\phi_v"]&\Gal(L^v/K_v).
\end{tikzcd}
\]
Combining these gives the commutative diagram
\[
\begin{tikzcd}[column sep=large,row sep=large]
\I_{K'}\arrow[r,"\phi_{L'/K'}"]\arrow[d,"\Nm"']&
\Gal(L'/K')\arrow[d,hook]\\
\I_K\arrow[r,"\phi_{L/K}"]&\Gal(L/K).
\end{tikzcd}
\]
The norm map on id\`eles sends $K'{}^\times$ into $K^\times$, so the lemma
follows from the earlier results.
\end{proof}
Example VII.8.2 and Lemma VII.8.4 show that part (a) holds for all
cyclotomic extensions of a number field $K$.
% END PRIMARY PROOF VII.8.1 BASE CHANGE

% BEGIN PRIMARY PROOF VII.8.1 CUP PRODUCT BRIDGE
\section*{The relation between the two statements}
\begin{authorresult}{Lemma}{8.5}
Let $L/K$ be an abelian extension of number fields. If part (b) holds for
$L/K$, then part (a) also holds. Conversely, if $L/K$ is cyclic and part
(a) holds, then part (b) also holds.
\end{authorresult}
\begin{proof}
Let $\chi\in\Hom(G,\Q/\Z)$. Regard $\chi$ as an element of
$H^1(G,\Q/\Z)$. Its image under the boundary map from
\[
 0\longrightarrow\Z\longrightarrow\Q\longrightarrow\Q/\Z
 \longrightarrow0
\]
is $\delta\chi\in H^2(G,\Z)$. Consider
\[
\begin{tikzcd}[column sep=large,row sep=large]
K^\times\arrow[r]\arrow[d,"\cup\delta\chi"']&
\I_K\arrow[r,"\phi_{L/K}"]\arrow[d,"\cup\delta\chi"]&
G\arrow[d,"\chi"]\\
H^2(G,L^\times)\arrow[r]&H^2(G,\I_L)\arrow[r]&\Q/\Z.
\end{tikzcd}
\]
The first two vertical arrows are cup product by $\delta\chi$.
If $\delta\chi$ is represented by the $2$-cocycle $n_{\sigma,\tau}$, then
the image of $x$ is represented by
$\sigma,\tau\longmapsto x^{n_{\sigma,\tau}}$.
The left square therefore commutes. The right vertical arrow is $\chi$
itself; commutativity of the right square follows from III.3.6.

Assume part (b) for $L/K$. Then
$\chi(\phi_{L/K}(a))=0$ for all characters $\chi$ of $G$, so
$\phi_{L/K}(a)=0$. Conversely, when $G$ is cyclic choose $\chi$ injective.
Then (a) implies (b), since $\cup\delta\chi$ is an isomorphism (II, 3.5).
\end{proof}
% END PRIMARY PROOF VII.8.1 CUP PRODUCT BRIDGE

% BEGIN PRIMARY PROOF VII.8.1 FINAL DESCENT
Part (a) is known for cyclotomic extensions, hence part (b) is known for
cyclic cyclotomic extensions. Proving part (b) proves the whole theorem.
The next result therefore completes its proof.
\begin{authorresult}{Lemma}{8.6}
If part (b) holds for cyclic cyclotomic extensions, then it holds for all
finite Galois extensions.
\end{authorresult}
\begin{proof}
Let $\beta\in H^2(/K)$. We are given that if $\beta\in H^2(L/K)$ for some
cyclic cyclotomic $L/K$, then $\sum_v\inv_v(\beta)=0$, where $\beta_v$ is
the image of $\beta$ in $H^2(K_v)$. Proposition VII.7.2 says that every
$\beta\in H^2(K)$ lies in $H^2(L/K)$ for such a cyclotomic extension.

The source appends the following bracketed explanation:
\begin{quote}
We use the inflation--restriction sequence
\[
 0\longrightarrow H^2(\Gal(L/K),L^\times)
 \longrightarrow H^2(G_K,\overline K^\times)
 \longrightarrow H^2(G_L,\overline L^\times),
\]
and $\beta$ is in the middle group and vanishes in the right-hand group.
\end{quote}
\end{proof}
% END PRIMARY PROOF VII.8.1 FINAL DESCENT

% BEGIN PRIMARY PROOF VII.8.1 RECIPROCITY QUOTIENT CONSEQUENCE
\section*{The author's reciprocity quotient consequence}
The author explains why part (a) gives the Reciprocity Law for the finite
abelian extension $L/K$ in that part.
It says that $K^\times$ is in the kernel of
$\phi_{L/K}:\I_K\to\Gal(L/K)$. The kernel already contains
$\Nm_{L/K}\I_L$, because this holds locally. Thus there is a homomorphism
\[
 \I_K/(K^\times\Nm_{L/K}\I_L)\longrightarrow\Gal(L/K).
\]
For a finite prime $v$ unramified in $L$, an id\`ele with a prime element
in the $v$-position maps to its Frobenius element. Proposition VII.4.7
therefore gives surjectivity. The second inequality (VII.5.1) gives
\[
 (\I_K:K^\times\Nm_{L/K}\I_L)\le[L:K],
\]
so this homomorphism is an isomorphism.
% END PRIMARY PROOF VII.8.1 RECIPROCITY QUOTIENT CONSEQUENCE

\medskip
\noindent\small
The single collected target is VII.8.1 with its two differently scoped
parts. The Brauer injectivity, cyclic cyclotomic splitting, local-action
context, cup-product bridge and quotient consequence are support for that
one target. VII.5.1 is linked prior support, not an additional count here.
The alternative route mentioned in Remark VII.8.3 is not added as a second
proof. This reader does not enlarge the target to the full
Brauer--Hasse--Noether exact sequence.
\end{document}
